
\section{Angular momentum: the big picture}
For a state \(\ket{j,m}\), the commuting operators \(\hat J^2\) and \(\hat J_z\) have
\[
\hat J^2\ket{j,m}=\hbar^2 j(j+1)\ket{j,m},\qquad
\hat J_z\ket{j,m}=\hbar m\ket{j,m}.
\]
The projections are \(m=-j,-j+1,\ldots,j\), giving \(2j+1\) magnetic substates.
\subsection{Adding two angular momenta}
If \(\hat{\mathbf J}=\hat{\mathbf J}_1+\hat{\mathbf J}_2\), then
\[
j=|j_1-j_2|,\ |j_1-j_2|+1,\ldots,j_1+j_2.
\]
The uncoupled basis \(\ket{j_1m_1}\ket{j_2m_2}\) and the coupled
basis \(\ket{(j_1j_2)jm}\) describe the same space:
\[
\ket{(j_1j_2)jm}
=\sum_{m_1+m_2=m}\langle j_1m_1,j_2m_2|jm\rangle
\ket{j_1m_1}\ket{j_2m_2}.
\]
\begin{methodbox}
\textbf{Exam method:} use the triangle rule to list allowed \(j\).
Check \(\sum_j(2j+1)=(2j_1+1)(2j_2+1)\).
\end{methodbox}
\subsection{Atomic quantum numbers and terms}
For one electron, \(\mathbf j=\mathbf l+\mathbf s\), where \(s=\tfrac12\).
For many-electron LS coupling,
\(\mathbf L=\sum_i\mathbf l_i\), \(\mathbf S=\sum_i\mathbf s_i\),
and \(\mathbf J=\mathbf L+\mathbf S\). A term is written
\[
{}^{2S+1}L_J,\qquad L=0,1,2,3,4\leftrightarrow S,P,D,F,G.
\]
The nuclear spin is \(I\); hyperfine angular momentum is
\(\mathbf F=\mathbf I+\mathbf J\). State which coupling scheme is valid.
